% TeX root = ../MathLogic.tex

	\chapter{Semantics notions for sets of formulas}

	Let $\Gamma \subseteq \fl$ ($\Gamma$ is a subset of the set of all formulas), then $\Gamma$ is called a \textbf{theory}.

	\begin{definition}
	  A \textbf{theory} $\Gamma \subseteq \fl$ is \textbf{satisfiable} \textbf{iff} $\exists v: \mathcal{L} \rightarrow \{0,1\}$ such that $\forall \gamma \in \Gamma$ then $\tilde{v} (\gamma) = 1$ (in symbols $v \models \Gamma$). \marginnote{$\leftarrow$ ($v \models \gamma$ we mean $v$ satisfies $\gamma$ as a formula, $\tilde{v} (\gamma) = 1)$}
	\end{definition}

	\begin{definition} A \textbf{theory} $\Gamma \subseteq \mathcal{F}_\mathcal{L}$ is \textbf{unsatisfiable} \textbf{iff} $\forall v: \mathcal{L} \rightarrow \{0,1\}$ there exists at least $\gamma \in \Gamma$ such that $\tilde{v} (\gamma) = 0$. (in symbols $v \not \models \Gamma$). \marginnote{$\leftarrow$ ($v \not \models \gamma$ we mean $v$ does not satisfy $\gamma$ as a formula, $\tilde{v} (\gamma) = 0$)}
	\end{definition}

	\section{Logical consequence}
	Of a formula $F \in \fl$ from a \textbf{theory} $\Gamma \subseteq \fl$.
		\begin{definition} Let $\Gamma \subseteq \FL$, $F \in \FL$ , then $F$ \textbf{is a logical consequence of} $\Gamma$, $\Gamma \models F$, \textbf{iff} for every assignment $\ta$ satisfying $\Gamma$ it holds that $v$ satisfies $F$ too:
	$$ v \models \Gamma \implies v \models F $$
	\end{definition}

	\noindent
	To express the fact that a given formula $F \in \FL$ \textbf{is a tautology} by means of the notion of logical consequence there are several ways, for example:
	$$ T \models F $$
	with $T$ being an always true constant. Another solution is to let $\Omega$ be a set of tautologies:
	$$\Omega \models F$$
	Or even:
	$$ \emptyset \models F$$
	since you can't falsify a formula in the $\emptyset$, there's no formula. For a \textbf{tautology} the standard notation most commonly used is:
	$$ \models F$$

	\subsection{Fundamental semantic properties of theories with respect to logical consequence}

	\begin{lemma}\label{lem:lemmunsat}
	Let $\Gamma \subseteq \FL$ be a theory and $A \in \FL$, then $\Gamma \models A$ \textbf{iff} $\Gamma \cup \{\neg A\}$ is unsatisfiable.
	\end{lemma}
	\begin{proof}
	  	\begin{align*}
		  \Gamma \models A &\textbf{ iff } \ \forall \ta, \ \tilde{v} (\gamma) = 1 \ \forall \gamma \in \Gamma \rightarrow \tilde{v} (A) = 1 & & \text{by definition of } \models \\
		  &\textbf{ iff } \ \forall \ta, \ \neg (\tilde{v} (\gamma) = 1 \ \forall \gamma \in \Gamma) \vee \tilde{v} (A) = 1 & & \text{by definition of } I_\rightarrow \\
		  &\textbf{ iff } \ \forall \ta, \ \neg ( \tilde{v} (\gamma) = 1 \ \forall \gamma \in \Gamma) \vee \tilde{v} (\neg A) = 0 & & \text{by definition of } I_\neg \\
		  &\textbf{ iff } \ \forall \ta, \ \exists \gamma \in \Gamma \ \tilde{v} (\gamma) = 0 \vee \tilde{v} (\neg A) = 0 & &  \\
		  &\textbf{ iff } \ \Gamma \cup \{\neg A\} \text{ is UNSAT } & & \text{by definition of unsatisfiable}
		\end{align*}
	\end{proof}

	\begin{lemma}[Semantical Deduction Lemma] Let $P,Q \in \FL$, $\{P\} \models Q$ \textbf{iff} $P \rightarrow Q$ is TAUTO.  The logical consequence relation can be internalized using the logical implication connective.
	\end{lemma}
	\begin{proof}
	  	  	\begin{align*}
		  \{P\} \models Q &\textbf{ iff } \ \forall \ta, \ \tilde{v} (P) = 1 \rightarrow \tilde{v} (Q) = 1 & & \text{by definition of } \models \\
		  &\textbf{ iff } \ \forall \ta, \ \tilde{v} (P \rightarrow Q) = 1 & & \text{by definition of } I_\rightarrow \\
		  &\textbf{ iff } \ P \rightarrow Q \text{ is TAUTO } & & \text{by definition of TAUTO}
		\end{align*}
	\end{proof}

	\begin{lemma}[Semantical deduction] Let $\Gamma \subseteq \FL$ and $P,Q \in \FL$, then $\Gamma \cup \{P\} \models Q$ (which will be also written as $\Gamma, P \models Q$) \textbf{iff} $\Gamma \models P \rightarrow Q$.
	\end{lemma}
	\begin{proof}
	  	\begin{align*}
		  \Gamma, P \models Q &\textbf{ iff } \ \forall \ta, \ (v \models \Gamma \wedge \tilde{v} (P) = 1) \rightarrow \tilde{v} (Q) = 1 & & \text{by definition of } \models \\
		  &\textbf{ iff } \ \forall \ta, \ v \models \Gamma \rightarrow \tilde{v} (P \rightarrow Q) = 1 & & \text{by definition of } I_\rightarrow \\
		  &\textbf{ iff } \ \Gamma \models P \rightarrow Q & & \text{by definition of } \models
		\end{align*}
	\end{proof}

	\section{Compactness Theorem}
	\begin{theorem}[Compactness Theorem] \hspace{1cm}
	\begin{center}
	  Let $\Gamma \subseteq \FL$, $\Gamma$ is SAT \textbf{iff} $\forall \ \Gamma' \subseteq_\omega \Gamma$,  $\Gamma'$ is SAT. \marginnote{$\leftarrow$ (with $\subseteq_{\omega}$ meaning finite subset.)}
	\end{center}
	\end{theorem} \marginnote{``IMPORTANT''}[-2cm]

	Do we \textit{really} need to consider the case $\Gamma$ is infinite (contains infinitely many formulas)? Yeah, it will be shown that in first order logic a single formula may contain the amount of information of infinitely many propositional formulas.

	Compactness theorem allows us to \textbf{reduce the analysis of satisfiability} of a possibly infinite theory \textbf{to the analysis of its finite sub-theories}, but observe that \textbf{there are infinitely many finite sub-theories of an infinite theory}.

	\begin{proof} \hspace{1cm}
	  \begin{itemize}
		\item ($\Rightarrow$) Proving the statement from \textbf{left to right} is trivial, if $v \models \Gamma$ then, clearly, by the definition of satisfiability and the notion of subset, $v \models \Gamma'$, for each $\Gamma' \subseteq \Gamma$ too (not even $\subseteq_\omega$, just any subset, even infinite).

		\item ($\Leftarrow$) Proving the statement from \textbf{right to left} is not so trivial.
	  \end{itemize}
	\end{proof}


	Consider the statement \textbf{equivalent} to the theorem \textbf{obtained via contraposition} ($A \rightarrow B$ says the same thing as $\neg B \rightarrow \neg A$):
	\begin{itemize}
		\item ($\Rightarrow$) If there is $\Gamma' \subseteq_\omega \Gamma$ UNSAT then clearly $\Gamma$ is UNSAT.

		\item ($\Leftarrow$) If $\Gamma$ is UNSAT then $\exists \Gamma' \subseteq_\omega \Gamma$ such that $\Gamma'$ is UNSAT
	  \end{itemize}

	\newpage

	\paragraph{Some terminology} before the proof: we say that a given theory $\Gamma \subseteq \FL$ is \textbf{finitely satisfiable} (FINSAT) \textbf{iff} $\forall \ \Gamma' \subseteq_\omega \Gamma$, $\Gamma'$ \textbf{is satisfiable}. \\

	\noindent
	We shall proceed with the proof of $\Gamma$ FINSAT $\implies$ $\Gamma$ SAT
	\begin{proof} \hspace{1cm} \vspace{-0.5cm}
	  \subsubsection{Preparation}
		We fix any \textbf{arbitrarily chosen listing} $F_1, F_2, \, ... \, , F_k , \, ...$ (enumerably infinite, without repetitions) of all formulas in $\FL$. To concretely implement such an enumeration we have two issues to face:

		\begin{itemize}
			 \item $1^{st}$ issue: the \textbf{alphabet of} $\FL$: $\{\wedge, \vee, \neg, \rightarrow\} \cup \{),(\} \cup \mathcal{L}$. This contains an infinite set an thus is an \textbf{infinite alphabet}. To \textbf{switch to a finite alphabet}:
			$$ \{\wedge, \vee, \neg, \rightarrow\} \cup \{),(\} \cup \{p,l\} $$

			with $p_k \in \mathcal{L}$ written as $p_k = plll \, ... \, l$, with $k$ many $l$s (just a unary encoding, we just need some kind of finite encoding, let it be a unary encoding, binary, base 10, hex, ...).\\


			\item $2^{nd}$ issue: take the alphabet $\{a,b, \, ... \, , \}$: list all the words on this alphabet, but they are infinite and we must define an order. If we use a lexicographic ordering we get:

			\begin{center}
				\begin{tabular}{| c |}
					\hline
					$\epsilon$ (empty word)\\
					$a$\\
					$aa$\\
					$a...aa$\\
					...\\
					$b \leftarrow$ never reached\\
					\hline
				\end{tabular}
			\end{center}

			So the letter $b$ appears after an infinite amount of words. Using a short-lex ordering might be fairer:
			\begin{center}
				\begin{tabular}{| c |}
					\hline
					$\epsilon$ (empty word)\\
					$a$\\
					... \\
					$z$ \\
					$aa$\\
					$ab$\\
					...\\
					$zz$\\
					\hline
				\end{tabular}
			\end{center}
			A ordering is \textbf{fair} if each word appears at a finite stage.
		\end{itemize}

		\newpage

		\subsubsection{Construction}

		We \textbf{define} an \textbf{infinite succession of theories} (sets of formulas) as follows:
		$$ D_0, D_1, \, ... \, , D_k , \, ... $$

		\noindent
		Inductively defined as follows:
		\begin{itemize}
			\item \textbf{Base:}
			$$D_0 := \Gamma$$

			\item \textbf{Step:}
			$$ D_{n+1} := \begin{cases}
				D_n \cup \{F_{n+1}\} \; & \text{ if } D_n \cup \{F_{n+1}\} \text{ is FINSAT} \\
				D_n \cup \{\neg F_{n+1}\}  & \text{ otherwise }
			\end{cases}
			$$
		\end{itemize}
		\textbf{N.B.} this definition is well given and in the second case does not directly imply that $D_n \cup \{\neg F_{n+1}\}$ is FINSAT, but just that $D_n \cup \{F_{n+1}\}$ is not FINSAT.

		 We also define:
		$$ D := \bigcup_{i \in \mathbb{N}} D_i$$

		Now we shall prove some facts about the sequence
		$$ \Gamma = D_0 \subseteq D_1 \subseteq D_2 \subseteq \, ... \, \subseteq D_k \subseteq \, ... \, \subseteq D $$

		To prove $\Gamma$ FINSAT $\implies$ $\Gamma$ SAT, we first need to prove some facts about the sequence above.

		\subsubsection{Fact 1} \label{fact:fact1}

		\begin{center}
		Every member of the sequence $D_n$ (for all $n \in \mathbb{N}$) is FINSAT.
		\end{center}

		\begin{proof}
			By induction on $n$
			\begin{itemize}
				\item \textbf{Base:} $n=0$, $D_0$ is FINSAT by definition, since $D_0 = \Gamma$ and $\Gamma$ is FINSAT by assumption.

				\item \textbf{Step:} Assume that Fact 1 is true for $D_0, D_1, \, ... \, , D_n$ and prove it for $D_{n+1}$. To prove this we shall proceed by contradiction. So we assume $D_{n+1}$ is not FINSAT with the aim of finding a contradiction. Absurdum Hypothesis: $D_{n+1}$ is not FINSAT.

				Then by the definition of the sequence, $D_n \cup \{F_{n+1}\}$ and $D_n \cup \{\neg F_{n+1}\}$ are both NOT FINSAT.

				Then by definition of FINSAT
				\begin{itemize}
					\item there exists
					$ D' \subseteq_\omega D_{n} \cup \{F_{n+1}\} $
					such that $D'$ is UNSAT.

					\item  there exists
					$ D'' \subseteq_\omega D_{n} \cup \{ \neg F_{n+1}\} $
					such that $D''$ is UNSAT.
				\end{itemize} \vspace{-0.25cm}

				\subsubsection{Lazy version}

				We can assume, without loss of generality that $F_{n+1} \in D'$, $\neg F_{n+1} \in D''$. We can argue that they are already in there, otherwise we can add them without major problems (if it's unsatisfiable adding something keeps it that way). \vspace{-0.25cm}

				\subsubsection{Actual version}

				If $F_{n+1} \notin D'$ then $D' \subseteq D_n$, but $D_n$ is FINSAT, so $D'$ is FINSAT, which is a contradiction.

				\newpage

				That is we can display $D'$ and $D''$ as follows:
				$$ D' = E' \cup \{F_{n+1}\}, \;\; D'' = E'' \cup \{\neg F_{n+1}\}$$
				with $E' = D' \setminus \{F_{n+1}\}$ and $E'' = D'' \setminus \{\neg F_{n+1}\}$.\\

				We observe that by construction that $E', E'' \subseteq_\omega D_n$.\\

				%End of L5

				\textbf{Note}: $E' \cup E'' \cup \{F_{n+1}\}$ is UNSAT, since it contains $D'$. For the same reason $E' \cup E'' \cup \{F_{n+1}\}$ is UNSAT since it contains $D''$. \\

				But $E' \cup E'' \finsat D_n$, so $E', E''$ \textbf{must be} FINSAT since we know that $D_n$ is FINSAT by I.H, so we get that $E' \cup E''$ is SAT.\\

				This \textbf{satisfiable} set \textbf{differs} from the \textbf{unsatisfiable} ones by \textbf{one formula}, notably the same formula.

				So, we have:
				\begin{itemize}
					\item $E' \cup E''$ \textbf{satisfiable}
					\item $E' \cup E'' \cup \{F_{n+1}\}$ \textbf{unsatisfiable} ($\dag$)
					\item $E' \cup E'' \cup \{\neg F_{n+1}\}$ \textbf{unsatisfiable} ($\dag \dag$)
				\end{itemize}

				By definition of satisfiability $\exists \ta$ such that $v \models E' \cup E''$ ($\tilde{v}(F) =  1$ for all $F \in E' \cup E''$).

				What about $\tilde{v}(F_{n+1})$?
				\begin{itemize}
					\item If it's $0$ it will satisfy $\tilde{v} (\neg F_{n+1})$ and then $v \models E' \cup E'' \cup \{\neg F_{n+1}\}$, \textbf{contradicting} ($\dag \dag$)
					\item If it's $1$ it will satisfy $\tilde{v} (F_{n+1})$ and then $v \models E' \cup E'' \cup \{F_{n+1}\}$, \textbf{contradicting} ($\dag$)
				\end{itemize}
				And, by \textbf{bivalence}, there are \textbf{no other cases to consider}. We have reached in any case the contradiction that proves the induction step ($D_{n+1}$ is actually FINSAT).
			\end{itemize}
		\end{proof}

		\subsubsection{Fact 2}

		$$ D = \bigcup_{i \in \mathbb{N}} D_i \;\; \text{ is FINSAT } $$

		Do we \textit{really} need to prove this? Yes, there are some properties of some sets, such as finiteness, that don't live into their union (with a collection of finite sets, their union might not be a finite set). \textbf{Properties of the members don't automatically carry over to their union}.

		\begin{proof}
			Let $D'$ be a generic finite subset of $D$ ($D' \subseteq_\omega D$). It is sufficient to prove that $D'$ is SAT.

			We can display the finite set $D'$ by listing its formulas:
			$$ D' = \{F_{i1}, F_{i2}, \, ... \, , F_{im}\} \; \text{for some } m \in \mathbb{N} $$

			We give these indexes since we don't know exactly what formulas are inside $D'$, but there must be a maximum, whatever that number is and whatever the numbers in the set are.\\

			Let $k$ be the maximum index
			$$ k = \max_{j=1}^m (ij) $$

			Then:
			\begin{align*}
			  D' &\subseteq D \\
			  D' &\subseteq \{F_1, \neg F_1, F_2, \neg F_2 \, ... \, , F_k, \neg F_k\}\\
			  D' &\subseteq \{F_1, \neg F_1, F_2, \neg F_2 \, ... \, , F_k, \neg F_k\} \cap D \subseteq D_k
			\end{align*}
			\marginnote{$\leftarrow$ (\textcolor{red}{CHANGED} from slides)}[-0.9cm]

			So $D' \subseteq_\omega D_k$ (as $D_k$ is defined in our sequence), but by Fact 1 $D_k$ is FINSAT. So $D'$ is satisfiable since it's one of its sub-theories.
			%End proof
		\end{proof}

		\subsubsection{Fact 3}

		\begin{center}
		For each formula $F_t \in \FL$ , \textbf{exactly one among} $F_t$ and $\neg F_t$ \textbf{belongs to} $D$.
		\end{center}
		\begin{proof}
			By construction of the sequence, $F_t \in D_t$ or $\neg F_t \in D_t$. Since $D_t \subseteq D$: $F_t \in D$ or $\neg F_t \in D$. This proves that at least one is in the set. \\

			We now have to exclude that both of them are in the set: ($\{F_t, \neg F_t\} \subseteq D$).\\
			But $D$ is FINSAT, and clearly $\{F_t, \neg F_t\}$ would contradict Fact 2 since it can't be satisfiable, by the definition of $I_\neg$. So $F_t \in D$ or $\neg F_t \in D$ and $\{F_t, \neg F_t\} \notin D$. Exactly one is $\in D$.
		\end{proof}

		\subsubsection{Construction}
		We define the following \textbf{assignment} $v_D: \mathcal{L} \rightarrow \{0,1\}$ such that $\forall p \in \mathcal{L}$, $v_D (p) = 1$ \textbf{iff} $p \in D$. 		\marginnote{
		$\leftarrow$ (Notice that $v_D$ is \textbf{well defined} (either $p \in D$ or $p \notin D$ for each possible $p \in \mathcal{L}$ also notice that $p \in \mathcal{L}$ its a letter so $p \in D$ means $p$ is inside one of the formulas of $D$)}

		\subsubsection{Fact 4}
		\begin{center}
		$v_D \models D$ \quad ($\forall F \in D$, $\tilde{v}_D (F) = 1$)
		\end{center}
		If we can prove this then we have essentially proved that $v_D \models \Gamma$ \textbf{since} $\Gamma \subseteq D$ (by construction of $D$). So $\Gamma$ is SAT.
		\begin{proof}
			By structural induction, we prove $\forall F \in \FL: \, \tilde{v}_D (F) = 1$ \textbf{iff} $F \in D$. \marginnote{$\leftarrow$ (Notice that would suffice to prove just one side of the iff, namely the side that says if $F \in D$ then $\tilde{v}_D (F) = 1$ (but the stronger proof is easier). We ask for the biconditional (iff)  for better using the I.H.)}

			\textbf{Base:} $F = p$, with $p \in \mathcal{L}$: so the statement we need to prove is the very definition of $v_D (p)$:
			$$\tilde{v}_D (p) = v_D (p) = 1 \; \textbf{ iff } \; p \in D $$

			\textbf{Step(s):} Proof by structural induction on the definition of formula and prove that it's $\in D$:
			\begin{itemize}
				\item $F = (\neg G)$, $F$ is the negation of some formula $G$.
				\begin{align*}
				  \tilde{v}_D (F) = 1 &\textbf{ iff } \tilde{v}_D (G) = 0 && \text{by definition of $I_\neg$} \\
				  &\textbf{ iff } G \notin D && \text{by I.H.} \\
				  &\textbf{ iff } \neg G \in D && \text{by Fact 3} \\
				  &\textbf{ iff } F \in D &&
				\end{align*}

				\item $F = (G \wedge H)$.
				  \begin{lemma*} $G \in D \wedge H \in D$ \textbf{ iff } $G \wedge H \in D$ \end{lemma*}
				  \begin{proof} \hspace{2cm}
					\begin{itemize}
					  \item ($\Rightarrow$) Lets suppose by absurdum $G \wedge H \notin D$
						\begin{align*}
						  G \wedge H \notin D &\textbf{ iff } \neg(G \wedge H) \in D && \text{by Fact 3}
						\end{align*}
						but then $\{G, H, \neg(G \wedge H)\} \subseteq_\omega D$ contraddicting Fact 2 since this set is obviously UNSAT \textcolor{red}{$\lightning$}
						\item ($\Leftarrow$) Lets suppose by absurdum $\neg(G \in D \wedge H \in D)$ that is $ \neg G \in D \vee \neg H \in D$ lets assume $\neg G \in D$ (its the same otherwise) \marginnote{$\leftarrow$ (destruct the assumption like in coq)} but then $\{G \wedge H, \neg G\} \subseteq_\omega D$ contraddicting Fact 2 since this set is obviously UNSAT \textcolor{red}{$\lightning$}
					\end{itemize}
				\end{proof}
				\vspace{-0.5cm}
				Now lets continue with the proof.
				\begin{align*}
				  \tilde{v}_D (F) = 1 &\textbf{ iff } \tilde{v}_D (G) = 1 \wedge \tilde{v}_D (G) = 1  && \text{by definition of $I_\wedge$} \\
				  &\textbf{ iff } G \in D \wedge H \in D  && \text{by I.H.} \\
				  &\textbf{ iff } G \wedge H \in D && \text{by Lemma} \\
				  &\textbf{ iff } F \in D &&
				\end{align*}

				\item $F = (G \vee H)$.
				  \begin{lemma*} $G \in D \vee H \in D$ \textbf{ iff } $G \vee H \in D$ \end{lemma*}
				  \begin{proof} \hspace{2cm}
					\begin{itemize}
					  \item ($\Rightarrow$) Lets suppose by absurdum $G \vee H \notin D$
						\begin{align*}
						  G \vee H \notin D &\textbf{ iff } \neg(G \vee H) \in D && \text{by Fact 3}
						\end{align*}
						but then assuming that $G \in D$ (its the same otherwise) \marginnote{$\leftarrow$ (destruct the assumption like in coq)} $\{G, \neg(G \vee H)\} \subseteq_\omega D$ contraddicting Fact 2 since this set is obviously UNSAT \textcolor{red}{$\lightning$}
						\item ($\Leftarrow$) Lets suppose by absurdum $\neg(G \in D \vee H \in D)$ that is $ \neg G \in D \wedge \neg H \in D$ but then $ \{\neg G, \neg H, G \vee H\} \subseteq_\omega D$ contraddicting Fact 2 since this set is obviously UNSAT \textcolor{red}{$\lightning$}
					\end{itemize}
				\end{proof}
				\vspace{-0.5cm}
				Now lets continue with the proof.
				\begin{align*}
				  \tilde{v}_D (F) = 1 &\textbf{ iff } \tilde{v}_D (G) = 1 \vee \tilde{v}_D (G) = 1  && \text{by definition of $I_\wedge$} \\
				  &\textbf{ iff } G \in D \vee H \in D  && \text{by I.H.} \\
				  &\textbf{ iff } G \vee H \in D && \text{by Lemma} \\
				  &\textbf{ iff } F \in D &&
				\end{align*}

				\item Prove implication for fun (remember that $G \rightarrow H$ has the same truth table as $(\neg G) \vee H$)
				%TODO
			\end{itemize}
		\end{proof}
		\marginnote{$\leftarrow$ (End of proof of Fact 4)}[-0.85cm]

		Since $v_D \models D$ then a fortiori (for stronger reasons) $v_D \models \Gamma$ (the assignment satisfies $D$, which is larger that $\Gamma$, so $v_D$ must satisfy $\Gamma$).

		So we have proved that:
		$$ \Gamma \; \text{ FINSAT } \Leftrightarrow \; \Gamma \; \text{ SAT } $$
	\end{proof}
	\marginnote{$\leftarrow$ (End of compactness proof)}[-0.85cm]
	%End of compactness proof

	\newpage

	\noindent
	By \textbf{contraposition}:
	$$ \Gamma \; \text{ UNSAT } \; \implies \; \exists \Gamma ' \subseteq_\omega \Gamma, \;\; \Gamma' \; \text{ UNSAT } $$

	\subsection{Remarks on Compactness Theorem:}

	$D$ is "Large" and satisfiable. We \textbf{cannot expand $D$ further without losing} the property of $D$ \textbf{being satisfiable}.

	\begin{proof}
		First of all, Since $D$ contains no contradiction its obviously not $\fl$, so $D \neq \FL$.

		Let $F_t \notin D$ then by Fact 3 $\neg F_t \in D$ and thus given $\overline{D} = D \cup \{F_t\}$ we have that $\{F_t, \neg F_t\} \subseteq_\omega \overline{D}$ so $\overline{D}$ is UNSAT.\\
		If I add anything not already in $D$, the negation of that formula must already be in $D$ so the set cannot remain satisfiable if I add anything else.\\
	\end{proof}
	\noindent
	$D$ is called a \textbf{maximal expansion}, with respect to \textbf{satisfiability}, of $\Gamma$. We call a set $E$ of formulas \textbf{maximally satisfiable} if it has the same property of $D$ that is $E$ is a set and for all $F \in \FL$, $F \notin E$ the set $E \cup \{F\}$ is UNSAT (if I add anything the set cannot remain satisfiable).
	$D$ is a maximal expansion of $\Gamma$, but there can be \textbf{many}, in our proof $D$ is fixed, after fixing the listing of formulas.\\

	% End L6

	If we reason about the \textbf{information contained} inside a \textbf{maximally satisfiable theory} $D$ we notice that it carries the \textbf{same amount of information} as a \textbf{truth value assignment} $v_D : \mathcal{L} \rightarrow \{0,1\}$ defined as $v_D (p) = 1$ \textbf{iff} $p \in D$ for all $p \in \mathcal{L}$. We can go the other way around, \textbf{starting from any assignment} $\ta$ and \textbf{construct the set of formulas} $D_v$:
	$$ D_v := \{F \in \FL : \, \tilde{v} (F) = 1\}$$

	\noindent
	This way \vspace{-0.15cm}
	\begin{itemize}
		\item $D_v$ is a MAXSAT theory
		\item $v_{D_{w}} = w$, with $w$ truth assignment $w: \rightarrow \{0,1\}$
		\item Same thing with: $D_{v_E} = E$ with $E$ a MAXSAT theory
	\end{itemize}

	\noindent
	Associated with every MAXSAT theory there is an assignment satisfying the theory ($v_D \models D$), and similarly for every assignment there can be a max sat theory satisfied by that assignment.\\

	\noindent
	How can we tell if:
	$$(\Gamma \subseteq \FL, \, A \in \FL) \;\;\;\;\; \Gamma \models A ?$$

	\noindent
	We have already proven in Lemma \ref{lem:lemmunsat} that we can \textbf{reduce this problem to} the unsatisfiability problem: $\Gamma \cup \{\neg A\}$ \textbf{is unsatisfiable}.

	\noindent
	Assume $\Gamma$ is a finite theory ($\Gamma$ contains finitely many formulas)
	$$\Gamma = \{B_1, B_2, \, ... \, , B_u \}$$ \vspace{-0.6cm}
	\begin{align*}
	\Gamma \models A &\textbf{ iff } & \\
	\{B_1, B_2, \, ... \, , B_u \} \models A &\textbf{ iff } & \\
	\{B_1, B_2, \, ... \, , B_u, \neg A \} \text{ is UNSAT (as a theory)} &\textbf{ iff } & \\
	B_1 \wedge B_2 \wedge \ldots \wedge  B_u \wedge \neg A \text{ is UNSAT (as a formula)} &\textbf{ iff } & \text{by definition of } I_\wedge
	\end{align*}

	\noindent
	What if $\Gamma$ is infinite? The conjunction of infinitely many things it's not a formula and its satisfiability it's not decidable in finite time.  But fortunately by compactness theorem, if $\Gamma  \cup \{\neg A\} \text{ is UNSAT}$ then there exists $ \Gamma' \finsat \Gamma \cup \{\neg A\} \text{ is UNSAT} $

Good news: we can have a finite theory which is unsatisfiable. Bad news: we don't know how to find it since there are infinitely many but this means that the procedure is not decidable but semidecidabile.

	\newpage

	\section{Semantic Equivalence (logical equivalence)}

	Are these the same?
	$$ A \wedge B \;\;\;\;\;\; B \wedge A $$
	These are syntactically different, they are different strings of symbols but semantically, \textbf{for all assignments} $\ta$ it holds that $\tilde{v} (A \wedge B) = \tilde{v} (B \wedge A)$, by the definition of $I_\wedge$. They always \textbf{get the same resulting truth values}.\\

	\subsection{Equivalence relation}
	%We need to introduce the equivalence relation: they're not the same formula, but the resulting value is always the same.\\

	\paragraph{What is a relation?} An $n$-ary relation $R$ on a set $S$ is just a subset of the $n$-th Cartesian power $S^n$ of $S$:
	$$ R \subseteq S^n$$
	A relation, mathematically, is just this.\\

	\begin{definition}[Equivalence relation]
	  	An equivalence relation $R$ on a set $S$, $R$ is a \textbf{binary relation} ($R \subseteq S^2$) such that it \textbf{satisfies 3 conditions}:
	\begin{itemize}
		\item $R$ is \textbf{reflexive}: $\forall s \in S$, $(s,s) \in R$
		\item $R$ is \textbf{symmetric}: $\forall s_1, s_2 \in S$ if $(s_1, s_2) \in R$ then also $(s_2, s_1) \in R$
		\item $R$ is \textbf{transitive}: $\forall s_1, s_2, s_3 \in S$ if $(s_1, s_2) \in R$ and $(s_2, s_3) \in R$ then $(s_1, s_3) \in R$
	\end{itemize}
	\end{definition}

	\begin{definition} Let $R \subseteq S^2$ be an equivalence relation, then for each $s \in S$ the 		\textbf{equivalence class} of $s$, called ($[s]_R$ or $\sfrac{s}{R}$) is \textbf{defined as}
	$$ [s]_R := \{t \in S: \, (s,t) \in R\} $$

	\noindent
	Observe: $(s,t) \in R$ \textbf{ iff } $[s]_R = [t]_R$.
	\end{definition}

	\begin{definition}[Partitions of a set] A partition of $S$ is a set $\{B_1, \, ... \, , B_k , \, ... \, \} = \{B_i : \, i \in  I\}$ of subsets of $S$ (not necessarily finite, $B_i \subseteq S$) such that:
	\begin{itemize}
		\item $B_i \cap B_j = \emptyset$  $\quad \forall i \neq j \in I$
		\item $\bigcup_{i \in I} B_i = S$
	\end{itemize}

	\noindent
	In simple words a partition is a way of \textbf{cutting $S$ into blocks}.
	\end{definition}

	\noindent
	Every \textbf{equivalence relation} $R \subseteq S^2$ \textbf{determines uniquely a partition} $P_R$ of $S$
	$$ P_R = \{[s]_R : \, s \in S \}$$

	\noindent
	We can go the other way around, every \textbf{partition} $P$ of $S$ \textbf{determines uniquely an equivalence relation} $R_P$ on $S$ ($R_P \subseteq S^2$)
	$$ R_p : (s,t) \in R_P \textbf{ iff } s,t \text{ belong to the unique same block } B_i \in P$$

	\subsection{Definition of Semantic equivalence}
	Let us verify the the following \textbf{is an equivalence relation on} $\FL$
	$$ \equiv \ \subseteq \FL^2$$

	\noindent
	For each $(A, B) \in \FL^2$ (or simpler, $A, B \in \FL$), then $(A, B) \in \ \equiv$, also written as $A \equiv B$ ($A$ is \textbf{logically equivalent to} $B$) \textbf{iff}, $\forall \ta$, $\tilde{v}(A) = \tilde{v} (B)$. Is it an equivalence relation?
	\begin{itemize}
		\item $\equiv$ is \textbf{reflexive}: for all $A \in \FL$, $A \equiv A$, for each $\ta$ $\tilde{v} (A) = \tilde{v}(A)$

		\item $\equiv$ is \textbf{symmetric}: for all $A,B \in \FL$, $A \equiv B$ implies $B \equiv A$, for all $\ta$: $\tilde{v}(A) = \tilde{v}(B)$ implies $\tilde{v}(A) = \tilde{v}(B)$

		\item $\equiv$ is \textbf{transitive}: for all $A,B,C \in \FL$, $A \equiv B$ and $B \equiv C$ implies $A \equiv C$, for all $\ta$: $\tilde{v}(A) = \tilde{v}(B)$ and $\tilde{v}(B) = \tilde{v}(C)$ implies $\tilde{v}(A) = \tilde{v}(C)$
	\end{itemize}

	\noindent
	$\equiv$ \textbf{determines a partition of} $\FL$. We subdivided the world in such a way that formulas that say the same thing are in the same block
	$$  \sfrac{\FL}{\equiv} = P_\equiv = \{[A]_\equiv : \, A \in \FL\}$$

	$[A \wedge B]_\equiv = [B \wedge A]_\equiv$ says that the formulas are equivalent, which is the same as saying $A \wedge B \equiv B \wedge A$.

	Two formulas are \textbf{logically equivalent} if they have the \textbf{same truth value under any assignment}. In general if $R \subseteq S^2$ is an equivalence relation, we write $\sfrac{S}{R} = \{[S]_R : s \in S\}$. This is called "$S$ modulo $R$" or "the quotient of $S$ over $R$".

	We eventually will enrich the set of logical equivalence classes $\sfrac{\FL}{\equiv}$ with operations which shall confer to the enriched set the status of an algebraic structure (namely: Boolean algebra). \vspace{-0.5cm}

	\section{Substitution (letter with a formula)}

	\textbf{Substitution} of a \textbf{propositional letter with a formula} (inside a formula). We want to get new formula from another one, while remaining semantically equivalent.

%	%Wtf, a cosa serve la prima parte
%	Intuitive idea: we have $(p \wedge \neg q)\left[t \rightarrow q/ p\right] = ((t \rightarrow q) \wedge q)$.
%	We substituted ...

	\noindent
	For each $A \in \FL$, for each $p \in \mathcal{L}$, for each $B \in \FL$, let us define the formula obtained by \textbf{replacing in} $A$ \textbf{each occurrence of} $p$ \textbf{with} $B$, $A\left[\sfrac{B}{p}\right] \in \FL$. Define it inductively on the definition of formula as follows:
	\begin{itemize}
		\item \textbf{Base:} $A = q$, $q \in \mathcal{L}$, define
		$$q[\sfrac{B}{p}] := \begin{cases}
			B \; & q = p  \\
			q \;  & q \neq p
		\end{cases}
		$$

		\item \textbf{Step:} cases, if we have to substitute:
		$$ (\neg A)[\sfrac{B}{p}] := \neg(A[\sfrac{B}{p}]) $$
		$$ (A_1 \wedge A_2) [\sfrac{B}{p}] := A_1 [\sfrac{B}{p}] \wedge A_2 [\sfrac{B}{p}] $$
		$$ (A_1 \vee A_2) [\sfrac{B}{p}] := A_1 [\sfrac{B}{p}] \vee A_2 [\sfrac{B}{p}] $$
		$$ (A_1 \rightarrow A_2) [\sfrac{B}{p}] := A_1 [\sfrac{B}{p}] \rightarrow A_2 [\sfrac{B}{p}] $$
	\end{itemize}

	\begin{lemma}[Substitution] \label{lemma:subst}
	Let $\ta$ be an assignment and $S,T \in \FL$. If $\tilde{v}(S) = \tilde{v} (T)$, then $\forall A \in \FL$, and $\forall p \in \mathcal{L}$
	$$ \tilde{v} (A [\sfrac{S}{p}]) = \tilde{v} (A[\sfrac{T}{p}]) $$
	\end{lemma}

	\begin{proof}
		By induction on $A$:
		\begin{itemize}
			\item \textbf{Base:} $A = q$, $q \in \mathcal{L}$, two cases:
			\begin{itemize}
				\item $q \neq p$, then $q [\sfrac{S}{p}] = q = q [\sfrac{T}{p}]$, then
				$$\tilde{v}(q[\sfrac{S}{p}]) = \tilde{v}(q) = \tilde{v} (q[\sfrac{T}{p}])$$

				\item $q = p$, then $q[\sfrac{S}{p}] = S$ and $q[\sfrac{T}{p}] = T$, so
				$$ \tilde{v} (q [\sfrac{S}{p}]) = \tilde{v} (S) = \tilde{v} (T) = \tilde{v} (q[\sfrac{T}{p}]) $$
			\end{itemize}

			\item \textbf{Step:} cases are:
			\begin{itemize}
				\item $A = \neg B$, by I.H. $\tilde{v} (B [\sfrac{S}{p}]) = \tilde{v} (B [\sfrac{T}{p}])$ we have that:
				$$ \tilde{v} ((\neg B) [\sfrac{S}{p}])  = I_\neg \left(\tilde{v} (B [\sfrac{S}{p}])\right) = I_\neg \left(\tilde{v} (B [\sfrac{T}{p}])\right) = \tilde{v} ((\neg B) [\sfrac{T}{p}])$$
				\item $A = (B \wedge C)$ , by I.H. $\tilde{v} (B [\sfrac{S}{p}]) = \tilde{v} (B [\sfrac{T}{p}])$ and $\tilde{v} (C [\sfrac{S}{p}]) = \tilde{v} (C [\sfrac{T}{p}])$ we have that:
				$$ \tilde{v} ((B \wedge C) [\sfrac{S}{p}]) = I_\wedge \left(\tilde{v} (B [\sfrac{S}{p}], \tilde{v} (C [\sfrac{S}{p}])\right) = I_\wedge \left(\tilde{v} (B [\sfrac{T}{p}], \tilde{v} (C [\sfrac{T}{p}])\right) = \tilde{v} ((B \wedge C) [\sfrac{T}{p}])$$
				\item $A = (B \vee C)$
				\item $A = (B \rightarrow C)$
			\end{itemize}
			The other cases are left as an exercise for the reader.
		\end{itemize}
	\end{proof}

	\begin{theorem}[Substitution]
	Let $S, T \in \FL$, if $S \equiv T$, then $\forall A \in \FL$ and $\forall p \in \mathcal{L}$:
	$$ A[\sfrac{S}{p}] \equiv A[\sfrac{T}{p}] $$
	\end{theorem}

	\begin{proof} \hspace{1cm} \\
		For all $\ta$ it holds that $\tilde{v} (s) = \tilde{v} (T)$, since $S \equiv T$. So by the Lemma \ref{lemma:subst} (substitution), we have $\tilde{v} (A [\sfrac{S}{p}]) = \tilde{v} (A[\sfrac{T}{p}])$, and then, by definition of $\equiv$
		$$ A[\sfrac{S}{p}] \equiv A[\sfrac{T}{p}] $$
	\end{proof}

	\begin{theorem}
	Given $F,G,H \in \FL$ the following are logical equivalences:
	\begin{itemize}
		\item \textbf{Idempotence}:
		  \begin{itemize}
			\item $F \wedge F \equiv F$
			\item $F \vee F \equiv F$
		  \end{itemize}
		\item \textbf{Commutativity}:
		  \begin{itemize}
			\item $F \vee G \equiv G \vee F$
			\item $F \wedge G \equiv G \wedge F$
		  \end{itemize}
		\item \textbf{Associativity}:
		  \begin{itemize}
			\item $F \wedge (G \wedge H) \equiv (F \wedge G) \wedge H$
			\item $F \vee (G \vee H) \equiv (F \vee G) \vee H$
		  \end{itemize}
		\item \textbf{Distributivity}:
		  \begin{itemize}
			\item $F \wedge (G \vee H) \equiv (F \wedge G) \vee (F \wedge H)$
			\item $F \vee (G \wedge H) \equiv (F \vee G) \wedge (F \vee H)$
		  \end{itemize}
		\item \textbf{Absorption}:
		  \begin{itemize}
			\item $F \vee (F \wedge G) \equiv F$
			\item $F \wedge (F \vee G) \equiv F$
		  \end{itemize}
	\end{itemize}
	\end{theorem}

	% End L7

	Exercise: prove that idempotence $F \vee F \equiv F$, $F \wedge F \equiv F$ follows only by absorption (using substitution theorem).

	\begin{itemize}
		\item \textbf{Double negation} (involutiveness of negation): $\neg \neg F \equiv F$

		\item \textbf{De Morgan}:
		\begin{itemize}
			\item $\neg (F \vee G) \equiv \neg F \wedge \neg G$
			\item $\neg (F \wedge G) \equiv \neg F \vee \neg G$
		\end{itemize}

		\item \textbf{Interdefinability}:
		\begin{itemize}
			\item $F \rightarrow G \equiv (\neg F) \vee G$
			\item $F \rightarrow G \equiv \neg (F \wedge \neg G)$
			\item $F \vee G \equiv (\neg F) \rightarrow G$
			\item $F \wedge G \equiv \neg (F \rightarrow \neg G)$
			\item $F \vee G \equiv \neg (\neg F \wedge \neg G)$
			\item $F \wedge G \equiv \neg (\neg F \vee \neg G)$
		\end{itemize}
	\end{itemize}

	\begin{proof}
		Just one case (the remaining ones are left as an exercise for the reader). We want to show that
		$$ \neg (F \vee G) \equiv \neg F \wedge \neg G$$
		For all $\ta$ $\tilde{v} (\neg (F \vee G)) = 1$
		\begin{itemize}[label*=]
			\item \textbf{iff} by the semantics of $\neg$: $I_\neg$: $\tilde{v} (F \vee G) = 0$
			\item \textbf{iff} by the semantics of $\wedge$: $\tilde{v} (F) = 0$ and $\tilde{v}(G) = 0$
			\item \textbf{iff} by the semantics of $\neg$: $\tilde{v} (\neg F) = 1$ and $\tilde{v} (\neg G) = 1$
			\item \textbf{iff} by the semantics of $\wedge$: $\tilde{v} (\neg F \wedge \neg G) = 1$
		\end{itemize}

		\noindent
		By bivalence $\neg (F \vee G) \equiv \neg F \vee \neg G$ (we've proved that it holds for 1, the only other possible value is 0 so we don't need to prove that).
	\end{proof} \vspace{-0.5cm}

	\section{Derived connectives}
	We define these operators as \textbf{shortened versions of longer formulas}.\\

	\noindent
	We define $F \leftrightarrow G$, as $(F \rightarrow G) \wedge (G \rightarrow F)$. Specify
	$$ I_\leftrightarrow: \, \{0,1\}^2 \rightarrow \{0,1\} $$
	\begin{center}
		\begin{tabular}{c c | c}
			$a$ & $b$ & $a \leftrightarrow b$ \\
			\hline
			0 & 0 & 1 \\
			0 & 1 & 0\\
			1 & 0 & 0\\
			1 & 1 & 1
		\end{tabular}
	\end{center}

	\paragraph{Observation:} $\models F \leftrightarrow G$ (is a tautology) \textbf{iff} $F \equiv G$ \textbf{iff} $\models \neg F \leftrightarrow \neg G$ \textbf{iff} $\models G \leftrightarrow F$ \textbf{iff} $\models \neg G \leftrightarrow \neg F$.\\

	\paragraph{Observation:} $F \models G$ and $G \models F$ \textbf{iff} $F \equiv G$.\\

	\noindent
	We define other derived connectives $\bot, \top$. Technically they are zero-ary connectives (and as such they are formulas in themselves), they're \textbf{logical constants}, don't need arguments to be complete to a formula.
	$$ \bot, \top \in \FL $$
	$$ I_\bot: \, \{0,1\}^0 \rightarrow \{0,1\}\;\; \text{ identify with } 0$$
	$$I_\top: \, \{0,1\}^0 \rightarrow \{0,1\} \;\; \text{ identify with } 1$$

	\noindent
	The $\{0,1\}^0$ represents the singleton set, so we can \textbf{identify the function with the result}. So
	$$ I_\bot := 0 \;\;\;\;\; I_\top := 1$$

	\noindent
	The $\top$ is in the \textbf{equivalence class of tautologies}:
	$$ \top \equiv A \textbf{ iff } A \text{ is an arbitrary tautology } (\models A), \text{ e.g. } \top \equiv p \rightarrow p$$

	\noindent
	And $\bot$ is in the \textbf{equivalence class of contradictions}:
	$$ \bot \equiv A \textbf{ iff } A \text{ is an arbitrary contradiction } (\models \neg A), \text{ e.g. } \bot \equiv p \wedge \neg p$$

	\noindent
	Some things \textbf{about these constants}:
	$$ \neg \top \equiv \bot, \;\;\; \neg \bot \equiv \top $$

	\noindent
	If $\models A$
	$$ \neg A \equiv A \rightarrow \bot $$
	$$ A \equiv \top \rightarrow A $$

	\noindent
	Some other equivalences:
	\begin{itemize}
		\item \textbf{Complementation}
		$$ A \wedge \neg A \equiv \bot $$
		$$ A \vee \neg A \equiv \top $$

		\item \textbf{Unit} and \textbf{neutral elements}
		$$ A \wedge \bot \equiv \bot $$
		$$ A \vee \bot \equiv A $$
		$$ A \wedge \top \equiv A $$
		$$ A \vee \top \equiv \top $$
	\end{itemize} \vspace{-1cm}

	\section{$\equiv$ as a congruence}

	$\equiv$ is \textbf{more than} simply being an \textbf{equivalence relation}. $\equiv$ is a \textbf{congruence} with respect to $\vee, \wedge, \rightarrow, \neg$ \textbf{thought as operations} as such:
	$$ \wedge : \, (\sfrac{\FL}{\equiv})^2 \rightarrow \sfrac{\FL}{\equiv} $$
	$$ ([A]_\equiv, [B]_\equiv) \rightarrow [A \wedge B]_\equiv $$
	If $A \equiv B$ and $C \equiv D$ then $A \wedge B \equiv C \wedge D$.

	\noindent
	Exercise: verify that $\equiv$ is a congruence with respect to all the connectives $\wedge, \vee, \rightarrow, \neg$ (also $\top, \bot, \leftrightarrow$ even if they're derived so there would technically be no need).

	\begin{proof}
		We demonstrate just the case of $\equiv$ being a congruence with respect to $\neg$.\\
		$A \equiv B$ \textbf{ iff } $\neg A \equiv \neg B$. $\forall \ta$, $\tilde{v}(A) = \tilde{v}(B)$ \textbf{ iff } $I_\neg (\tilde{v}(A)) = I_\neg (\tilde{v}(B))$ \textbf{ iff } $\tilde{v}(\neg A) = \tilde{v}(\neg B)$.
	\end{proof}

	\section{Partial order relation}

	Let $R$ be a \textbf{binary relation} on $S$, $R \subseteq S^2$ such that:
	\begin{itemize}
		\item $R$ is reflexive $(s, s) \in R$  $\ \forall s \in S$
		\item $R$ is \textbf{antisymmetric} $\ \forall (s,t) \in R$, if $(s,t) \in R$ and $(t,s) \in R$ then $s = t$
		  \marginnote{$\leftarrow$ (don't confuse this property with the asymmetric property)}
		\item $R$ is transitive if $(s,t) \in R$ and $(t,r) \in R$ then $(s,r) \in R$
	\end{itemize}
	Any relation with these properties is a partial order relation (poset). \marginnote{$\leftarrow$ (Symbols: $\leq, \preceq, \sqsubseteq$.)}

	\noindent
	Some examples of partial orderings for the naturals:
	$$ (\mathbb{N}, \leq) $$
	$a \leq b$ if it exists $c \in \mathbb{N}$ such that $a + b = c$. This relation has a minimum element, which is $0$.

	$$ (\mathbb{N}, \geq) $$
	$a \geq b$ if it exists $c \in \mathbb{N}$ such that $a - b = c$. This relation has a maximum element, which is $0$. Both of the previous posets are also \textbf{total orders}, since this property holds:
	\begin{itemize}
		\item $R$ is \textbf{total} $\forall a,b \in S \ (a, b) \in R$ or $(b, a) \in R$
	\end{itemize}

	\noindent
	Another unituitive poset is this one:
	$$ (\mathbb{N}, =) $$
	In this relation every element is on the same level.\\

	\noindent
	Divisibility order:
	$$ (\mathbb{N}, |) $$
	$a|b$ ($a$ divides $b$) if there is $c \in \mathbb{N}$ such that $a \cdot c = b$.\\

	\begin{center}
			\begin{tikzpicture}[scale=0.7]
			\node (1) at (2,0) {1};
			\node (2) at (-2,2) {2};
			\node (3) at (2,2) {3};
			\node (4) at (-2,4) {4};
			\node (6) at (2,4) {6};
			\node (0a) at (2,5) {...};
			\node (1b) at (6,2) {...};
			\node (0b) at (-2,5) {...};
			\node (0c) at (6,5) {...};
			\node (0) at (2,7) {0};

			\draw (0) -- (0a);
			\draw (0) -- (0b);
			\draw (0) -- (0c);
			\draw (1) -- (2);
			\draw (1) -- (3);
			\draw (1) -- (1b);
			\draw (2) -- (4);
			\draw (2) -- (6);
			\draw (3) -- (6);
		\end{tikzpicture}
	\end{center}

	This order has $1$ as a minimum, all the prime numbers in the first level (atoms), then all numbers created just by primes and so on until the maximal element, which is zero since there's always a number such that $n \cdot 0 = 0$.

	A \textbf{set} $P$, equipped with a partial order relation $\leq \ \subseteq P^2$ is called a \textbf{partially ordered set}, for short \textbf{poset} $(P, \leq)$.\\

	\section{Lattice}
	a \textbf{Lattice} (Reticolo) is a poset such that each pair of element in the poset has a unique \textbf{greatest lower bound} and a unique \textbf{lowest upper bound}. where considering $S \subseteq P$, $P$ poset, $(P, \leq)$.\\

	\noindent
	We can define the \textbf{infimum} (greatest lower bound): $\inf S$ is the element $t \in P$ such that $\forall s \in S$ then $t \leq s$ and for all $z \in P$, if $z \leq s$ for all $s \in S$ then $z \leq t$ ($t$ is the greatest lower bound). \marginnote{$\leftarrow$ (the infimum for a given set can also not exists)}\\

	\noindent
	We can define the \textbf{supremum} (lowest upper bound) $\sup S$: is the element $t \in P$ such that $\forall s \in S$ then $s \leq t$ and for all $z \in P$, if $s \leq z$ for all $s \in S$ then $t \leq z$ ($t$ is the lowest upper bound).\marginnote{$\leftarrow$ (the supremum for a given set can also not exists)}

	\noindent
	We can also define:
	$$ \min S = t \textbf{ if } t = \inf S \text{ and } t \in S$$
	$$ \max S = t \textbf{ if } t = \sup S \text{ and } t \in S$$

	\begin{example*}
	of a set lacking $\sup$ (and then lacking $\max$)
	$$ S \subseteq \mathbb{Q} \;\;\;\;\; S := \{x : x^2 \leq 2\}$$
	This has a $\sup$, but not in the rationals. It would be $\sqrt{2} \in \mathbb{R} \setminus \mathbb{Q}$.\\

	\noindent
	While
	$$ S' \subseteq \mathbb{Q} \;\;\;\;\; S' := \{x : x \lneq 1\}$$
	Has a $\sup$ but not a $\max$, $\sup S' = 1$ but $1 \notin S'$, so $\max S'$ does not exist.

	$$ S'' \subseteq \mathbb{Q} \;\;\;\;\; S'' := \{x : x \leq 1\}$$
	$\sup S'' = 1$, $1 \in S''$ so $\max S'' = 1$
	\end{example*}

	\noindent
	Exercise: is our poset a lattice?
	$$ (\mathbb{N}, \leq) $$
	it's a lattice because every pair has an greatest lower bound (the smaller number) and a lowest upper bound (the largest number).

	$$ (\mathbb{N}, |) $$
	Divisibility is a lattice, the inf is the greatest common divisor, while the sup is the lowest common 	multiple.

	$$ (\mathbb{N}, =) $$
	Is NOT a lattice, every element is on the same level, there is no $\inf$ or $\sup$.

	%End of  L8

%	Esercizio random?
%
%	Prove idempotence, follows from absorption
%	%Check and try to prove this from notes
%
%
%	Using explicitly the substitution theorem:
%
%	Take p \in L not occurring in $F$
%	$$ (F \wedge p) [F/p] = F \wedge F \equiv (F \wedge p)[F \wedge (F \vee F)/p] = F \vee (F \wedge (F \vee F)) \equiv F$$
%
%	By absorption and substitution theorem we can say that $F \vee F \equiv F$.\\
%
%	***

	\section{Algebraic Structure}

	%Lattices are not only posets with an order structure, but they're also an algebraic structure, each boolean algebra can be considered ...

	\paragraph{(Very sketchy) Definition of Algebraic Structure:} Let $S$ be a set, and $\ast_1, \ast_2, \, ... \, , \ast_n$ a number of operations (can be infinite) on $S$ (that is functions $\ast_i : S^{\nu(i)} \rightarrow S$, where $\nu(i)$ gives the arity of the operation). \\

	\noindent
	Then $(S, \ast_1, \ast_2, \, ... \, , \ast_n)$ is an \textbf{algebraic structure}, or is an algebra. \\

	\begin{example}
	  $$ (\mathbb{Z}, +, -, 0) $$
	  \vspace{-1cm}
	  \begin{align*}
	+ &: \mathbb{Z}^2 \rightarrow \mathbb{Z}  &(a,b) \rightarrow a + b  \\
	- &: \mathbb{Z}^1 \rightarrow \mathbb{Z}  &a \rightarrow - a \\
	0 &: \mathbb{Z}^0 \rightarrow \mathbb{Z}  &0 \text{ the constant/integer number } 0
	\end{align*}
	The arity of these operations is, respectively: $\nu(+) = 2$, $\nu(-) = 1$, $\nu(0) = 0$.
	\end{example}

	\subsection{Isomorphism}
	Two \textbf{structures} are \textbf{isomorphic}
	$$(S, \ast_1, \ast_2, \, ... \, , \ast_n) \cong (R, \odot_1, \odot_2, \, ... \, , \odot_m) $$

	\noindent
	If they mathematically behave in the same manner, i.e. \textbf{iff}:
	$$ m = n $$
	$$ \nu_S^{(i)} = \nu_R^{(i)} \;\; \text{ for each } i = 1,2, \, ... \, , n $$
	$$ \exists f : S \rightarrow R \text{ such as: } $$
	\vspace{-0.6cm}
	\begin{align*}
		\begin{cases}
		f : \text{ injective: }  &x \neq y \in S \implies f(x) \neq f(y) \\
		f : \text{ surjective: }  &\forall y \in R \text{ there is } x \in S \text{ s.th. } f(x) = y
	\end{cases}
	\end{align*}
	$$\implies f \text{ is bijective (1-1 correspondance)} $$
	$$ f: \text{ is a homomorphism}: f(\ast_i (s_1, \, ... \, , s_k)) = \odot_i(f(s_1), f(s_2), \, ... \, , f(s_k))$$
	$$ \nu_S^{(i)} = k, \text{ for all } (s_1, \, ... \, s_k) \in S^k $$

	\marginnote{$\leftarrow$ (guarda slide in piu)}
	\vspace{-0.15cm}

	\subsection{Lattices as algebraic structures}

	A \textbf{lattice as an algebraic structure} is a \textbf{set with binary operations}
	$$ (R, \sqcap, \sqcup) \;\;\;\;\;
	\begin{array}{c}
		\sqcap: R^2 \rightarrow R\\
		\sqcup: R^2 \rightarrow R
	\end{array}
	 $$

	\noindent
	such that \textbf{both} these \textbf{operations} are \textbf{associative and commutative}.
	$\sqcap, \sqcup$ are related by absorption:
	$$
	\forall\; x, y \in R \;\;\;\;\;
	\begin{cases}
		x \sqcap (x \sqcup y) = x \\
		x \sqcup (x \sqcap y) = x
	\end{cases}$$

	\noindent
	A \textbf{lattice as a poset} can be \textbf{redefined as/associated} to a \textbf{lattice as an algebraic structure}
	$$ (R, \leq ) \stackrel{unique}{\leftrightsquigarrow} (R, \sqcap, \sqcup)$$
	each lattice meant as a poset can be thought of uniquely as an algebra and vice versa.
	from left to right:
	$$ (R, \leq ) \rightsquigarrow (R, \sqcap, \sqcup) $$
	we just have to define the operations. For all $a,b \in R$ we need to define:
	$$ a \sqcap b := \inf \{a, b\}$$
	$$ a \sqcup b := \sup \{a,b\}$$

	\newpage

	While from right to left
	$$ (R, \sqcap, \sqcup ) \rightsquigarrow (R, \leq) $$

	For all $a,b \in R$, $a \leq b$ \textbf{iff} $a = a\sqcap b$ (or, equivalently, $b = a \sqcup b$).

	Check: $a \leq b$ defined as above is such that $\leq$ is:
	\begin{itemize}
		\item reflexive: for all $a \in R$, $a \leq a$ iff $a = a \sqcap a$ (by idempotence)
		\item antisymmetric: ...
		\item transitive: ...
	\end{itemize}

	We now proved that $(R, \leq)$ is a poset, but the lattice is a special kind of poset. So now we need to check that the poset $(R, \leq)$ is indeed a lattice. We need to prove that \textbf{for each} $a,b \in R$ \textbf{there is} $\inf\{a,b\}$ \textbf{and} $\sup \{a,b\}$.

	Letting $c \leq a$ and $c \leq b$ we have to show that $c \leq a \sqcap b$, since it means we've proven that $a \sqcap b = \inf\{a,b\}$, so it exists and is defined, $\inf\{a,b\} = a \sqcap b$.\\
	$$ c \underbrace{=}_{\text{def of } c \leq a} = a \sqcap c \underbrace{=}_{\text{def of } c \leq b} a \sqcap (b \sqcap c) \underbrace{=}_{\text{assoc}} (a \sqcap b) \sqcap c \; \underbrace{\text{ iff }}_{\text{def of } \leq} \; c \leq a \sqcap b$$
	The proof that $\sup\{a,b\} = a \sqcup b$ is analogous. We have proved that $(R, \leq)$ is a lattice (as a poset).

	\section{Boolean Algebras}

	A Boolean Algebra is a particular type of algebra:
	$$ (S, \sqcap, \sqcup, ()^c, 0, 1 ) $$

	\noindent
	The arity of the operations is:
	\begin{itemize}
		\item $\sqcap, \sqcup : S^2 \rightarrow S$
		\item $()^c: S \rightarrow S$ \marginnote{$\leftarrow$ (The $()^c$ is the complement)}
		\item $0,1: S^0 \rightarrow S$
	\end{itemize}

	\noindent
	\textbf{Additional properties}: $(S, \sqcap, \sqcup)$ is a lattice (has an order), which is:
	\begin{itemize}
		\item \textbf{Bounded:} for all $a \in S$ there's a minimum $0 \leq a$ and a maximum $a \leq 1$ such that
		$$ \begin{cases}
			a \sqcap 0 = 0 \;\; &  a \sqcup 1  = 1 \\
			a \sqcup 0 = a & a \sqcap 1 = a
		\end{cases}$$

		\item \textbf{Complemented:} for all $a \in S$, such that $a \sqcap a^c = 0$ and $a \sqcup a^c = 1$.\\

		\item \textbf{Distributive:} for all $a,b,c \in S$, $a \sqcap (b \sqcup c)$
		$$ \begin{array}{c}
			a \sqcap (b \sqcup c) = (a \sqcap b) \sqcup (a \sqcap c) \\
			a \sqcup (b \sqcap c) = (a \sqcup b) \sqcap (a \sqcup c)
		\end{array}$$
	\end{itemize}

	So, boolean algebra is a shorthand name for bounded, complemented and distributive lattices.\\

	\newpage

	\begin{example} Let $A$ be a set and let $\mathcal{P}(A)$ be the powerset of $A$ (set of all subsets of $A$):
	$$ (\mathcal{P}(A), \cap, \cup, \overline{()}, \emptyset, A) $$

	\noindent
	Check that this is a boolean algebra. It must be commutative, absorption, bounds, ...\\
	If $A$ is finite, say $|A| = n$, then $|\mathcal{P}(A)| = 2^n$. $2^A$ is a commonly used notation for the powerset.
	\end{example}
	\noindent
	Every \textbf{finite} Boolean algebra has $2^k$ elements for some $k \in \mathbb{N}$.\\
	Every \textbf{finite} Boolean algebra can be conceived/is \textbf{isomorphic with} the \textbf{powerset of some finite set}.

	$(\mathcal{P}(A), \subseteq)$, forall $S,T$, $S \subseteq T$ \textbf{iff} $S = S \cap T$ (or $T = S \cup T$). This is the lattice isomorphic to the boolean algebra defined above.

	\begin{example}
	$$A = \{a,b\}, \; |\mathcal{P} (A)| = |2^A| = 2^2 = 4$$
	$$ 2^A = \{\emptyset, \{a\}, \{b\}, \{a,b\}\} $$
	So
	$$ (2^A, \subseteq)$$
	is a lattice.
	\end{example}

	\subsection{Boolean algebra of truth-values:}
	$$ (\{0,1\}, I_\wedge = \min, I_\vee = \max, I_\neg = 1-x, 0 = I_\bot, 1 = I_\top) $$
	This is the smallest significant boolean algebra and in the ordering we have that 0 is below 1.
	\noindent
	It's isomorphic to:
	$$ \cong (\mathcal{P}(\{\ast\}), \cap, \cup, \overline{()}, \emptyset, \{\ast\}) $$ \marginnote{$\leftarrow$ ($\{*\}$ is the singleton set)}[-0.8cm]

	\vspace{-0.4cm}

	\noindent
	Boolean algebras of $n$-tuples of truth-values:
	$$ (\{0,1\}^n, \overline{\min}, \overline{\max}, \overline{1} - x, (0,0, \, .... \, , 0), (1,1, \, ... \, , 1)) $$

	\noindent
	Operations are defined pointwise (componentwise):
	$$ \overline{\min} ((b_1, \, ... \, , b_n), (c_1, \, ... \, c_n)) := (\min \{b_1, c_1\}, \, ... \, , \min \{b_n, c_n\})$$
	$$ \overline{\max} ((b_1, \, ... \, , b_n), (c_1, \, ... \, c_n)) := (\max \{b_1, c_1\}, \, ... \, , \max \{b_n, c_n\})$$
	$$ \overline{1} - (b_1, \, ... \, , b_n) = (1-b_1, \, ... \, , 1-b_n) $$

	\subsection{Lindembaum Algebras of (equivalence classes of) formulas}
	It is a boolean algebra.
	$$ (\FLe, \bm{\wedge}, \bm{\vee}, \bm{\neg}, [\bot]_\equiv, [\top]_\equiv) $$

	Since $\equiv$ is a congruence with respect to $\wedge, \vee, \neg$ the following \textbf{definition is well given} (does not depend on the representatives of the classes). For all $A, B \in \FL$:
	$$ [A]_\equiv \bm{\wedge} [B]_\equiv := [A \wedge B]_\equiv $$
	$$ [A]_\equiv \bm{\vee} [B]_\equiv := [A \vee B]_\equiv $$
	$$ \neg [A]_\equiv := [\neg A]_\equiv $$

	\noindent
	It has an \textbf{order relation} since it's also a \textbf{lattice}:
	$$ (\FLe, \leq) \;\;\;\;\; [A]_\equiv \leq [B]_\equiv  \; \text{ iff } \; [A]_\equiv = [A \wedge B]_\equiv \text{ or } [B]_\equiv = [A \vee B]_\equiv $$

	\noindent
	Some combinatorics about $\FLe$
	$$ |\FL| = |\mathcal{L}| = |\mathbb{N}| = \aleph_0 $$
	$$ |\FLe| = \aleph_0 $$
	No powerset has cardinality $\aleph_0$, so $|\FLe|$ is \textbf{not isomorphic} with any powerset algebra.

	By $\FL^{(n)}$ we mean the set of all formulas built using only the first $n$ propositional letters $\{p_1, p_2, \, ... \, , p_n\}$. Check that  %Complete
	$$ |\FL^{(n)}/\equiv|$$
	is a boolean algebra (called Lindenbaum algebra over the first $n$ letters).\\

	\begin{figure}[H]
	\centering
	\includegraphics[width=1\textwidth]{img1.png}
	\end{figure}

	%End L9

	\newpage

	\noindent
	First exercise: taking
	$$ (\FLe, \leq)$$
	thought as a poset, we know that
	$$ [A]_\equiv \leq [B]_\equiv $$

	\begin{itemize}[label*=]
		\item \textbf{iff} $[A]_\equiv = [A \wedge B]_\equiv $
		\item \textbf{iff} $A \equiv A \wedge B$
		\item \textbf{iff} $\models A \leftrightarrow (A \wedge B)$
		\item \textbf{iff} $\models A \rightarrow B$
	\end{itemize}

	Looking at the truth table, we can see that it's structured in such a way that no assignment can make $A \rightarrow B$ a tautology.\\

	\noindent
	Exercise: $F^{(n)}/\equiv$ is isomorphic with
	$$2^{2^{n}} \cong (\mathcal{P} (\mathcal{P} (\, ... \, )), \cap, \cup, ()^c, \emptyset, \, ...) $$

	\paragraph{Some other notable logically equivalent formulas:}
	\begin{itemize}
		\item \textbf{De Morgan Laws on $n$ many variables}:
		$$ \neg (F_1 \vee F_2 \vee \, ... \, , \vee F_n) \equiv \neg F_1 \wedge \neg F_2 \wedge \, ... \, \wedge \neg F_n $$
		$$ \neg (F_1 \wedge F_2 \wedge \, ... \, , \wedge F_n) \equiv \neg F_1 \vee \neg F_2 \vee \, ... \, \vee \neg F_n $$
		\item \textbf{Generalized distributivity}:
		$$ (F_1 \vee \, ... \, \vee F_u) \wedge (G_1 \vee \, ... \, \vee G_v) \equiv \bigvee_{i=1}^u \bigvee_{j = 1}^v (F_i \wedge G_j) $$
		$$ (F_1 \wedge \, ... \, \wedge F_u) \vee (G_1 \wedge \, ... \, \wedge G_v) \equiv \bigwedge_{i=1}^u \bigwedge_{j = 1}^v (F_i \vee G_j) $$
	\end{itemize}

	\section{Term Function}
	%Motivations for a new way to look at things.\\
	Motivations: $\models F$ \textbf{iff} (by definition) for all $\ta$ it holds that $\tilde{v}(F) = 1$.

	\noindent
	We Desire that:
	\begin{itemize}
		\item the \textbf{formula} "becomes" the \textbf{function}
		\item the \textbf{assignments} "become" the \textbf{arguments}
	\end{itemize}

	We fix a \textbf{formula} $A \in \FL$ and let $Var(A)$ be the set of \textbf{propositional letters actually occurring} in $A$ (since assignments are infinite by definition). Now fix a \textbf{suitable natural number} $n \in \mathbb{N}$ such that $Var(A) \subseteq \{p_1, p_2, \, ... \, , p_n\}$ (just take $n$ "large enough" for $A$).\\

	\textbf{Desideratum} (what we want): we want to \textbf{define a function} $\hat{A}: \{0,1\}^n \rightarrow \{0,1\}$ \textbf{such that} for all assignments $\ta$:
	$$ \hat{A} (v(p_1), v(p_2), \, ... \, , v(p_n)) = \tilde{v} (A) $$
	The \textbf{formula} has \textbf{become a function}, the \textbf{assignment} has become the \textbf{argument}. Sometimes it's more convenient.\\

	\newpage

	\textbf{Definition}: of \textbf{term function} $\hat{A}$ associated with the formula $A \in \FL$ with respect to $n \in \mathbb{N}$ ($Var(A) \subseteq \{p_1, \, ... \, , p_n\}$).\\

	We define this by \textbf{induction on the structure of} $A$:
	\begin{itemize}
		\item \textbf{Base}: $A = p_i$, $p_i \in \{p_1, \, ... \, , p_n\}$
		$$ \hat{A}: \{0,1\}^n \rightarrow \{0,1\}$$
		$$ \hat{p}_i: \{0,1\}^n \rightarrow \{0,1\}$$

		for all $(t_1, \, ... \, , t_n) \in \{0,1\}^n$
		$$ \hat{p}_i (t_1, \, ... \, , t_n) :=  t_i $$

		so it satisfies the desideratum:
		$$ \hat{p}_i (v(p_1), \, ... \, , v(p_i), \, ... \, , v(p_n)) = v(p_i) = \tilde{v} (p_i) $$

		\item \textbf{Steps}: we have to define the cases:
		\begin{itemize}
			\item $A = \neg B$
			$$ \hat{A} (t_1, \, ... \, , t_n) = \hat{(\neg B)} (t_1, \, ... \, , t_n) := I_\neg (\hat B (t_1, \, ... \, , t_n)) =  1 - (\hat B (t_1, \, ... \, , t_n))$$

			\item $A = B \wedge C$
			$$ \hat{A} (t_1, \, ... \, , t_n) =  \hat{(B \wedge C)} (t_1, \, ... \, , t_n) := I_\wedge (\hat B (t_1, \, ... \, , t_n), \hat C (t_1, \, ... \, , t_n))$$

			\item $A = B \vee C$
			$$ \hat{A} (t_1, \, ... \, , t_n) =  \hat{(B \vee C)} (t_1, \, ... \, , t_n) := I_\vee (\hat B (t_1, \, ... \, , t_n), \hat C (t_1, \, ... \, , t_n))$$

			\item $A = B \rightarrow C$
			$$ \hat{A} (t_1, \, ... \, , t_n) =  \hat{(B \rightarrow C)} (t_1, \, ... \, , t_n) := I_\rightarrow (\hat B (t_1, \, ... \, , t_n), \hat C (t_1, \, ... \, , t_n))$$

			\item $\top$
			$$ \hat{\top} (t_1, \, ... \, , t_n) = 1$$

			\item $\bot$
			$$ \hat{\bot} (t_1, \, ... \, , t_n) = 0$$
		\end{itemize}
	\end{itemize}

	\subsection{Boolean algebra of term functions}

	$$ \hfln := \left( \left\{\hat F: \{0,1\}^n \rightarrow \{0,1\}: F \in \fln \right\}, \hat \wedge, \hat \vee, \hat \neg, \hat \top, \hat \bot \right)$$
	where
	$$ (\hat F \wedge \hat G) (t_1, \, ... \, , t_n) = I_\wedge \{\hat F (t_1, \, ... \, , t_n), \hat G (t_1, \, ... \, , t_n)\}$$
	$$ (\hat F \vee \hat G) (t_1, \, ... \, , t_n) = I_\vee \{\hat F (t_1, \, ... \, , t_n), \hat G (t_1, \, ... \, , t_n)\}$$
	$$ \hat \neg \hat F (t_1, \, ... \, , t_n) = I_\neg (\hat F (t_1, \, ... \, , t_n)) $$

	We can prove that $\hfln$ is \textbf{isomorphic} with $\flne$ ($\hfln \cong \flne$) via: $\hat F \mapsto [F]_\equiv$ for each $F \in \FL$.

	$$\hfln = \left(\{ \hat F: \{0,1\}^n \rightarrow \{0,1\}: F \in \FL\}, \hat \wedge, \hat \vee, \hat \neg, \hat \top, \hat \bot \right)$$

	\noindent
	Compare with
	$$ Func^{(n)} = \left(\{ g: \{0,1\}^n \rightarrow \{0,1\} \}, \hat \wedge, \hat \vee, \hat \neg, \hat \top, \hat \bot \right)$$

	\noindent
	Check that $Func^{(n)}$ \textbf{is a boolean algebra}.\\

	\noindent
	This
	$$ \fln \subseteq Func^{(n)}$$
	is trivially true since each term function belongs to $Func^{(n)}$, while
	$$ \fln \supseteq Func^{(n)}$$
	is also true, called \textbf{functional completeness of propositional logic}.\\

	\noindent
	So
	$$ \fln = Func^{(n)} $$
	They're not isomorphic, it's \textbf{literally the same algebra}, they have the same universe and same operations.

	\section{Functional Completeness Theorem}

	\begin{theorem}
	We have to prove that \textbf{for each natural number} $n \in \mathbb{N}$, the set of function $Func^{(n)} \subseteq \hfln$. That is, for each function $g: \{0,1\}^n \rightarrow \{0,1\}$ there exists a formula $F \in \fln$ such that  $g = \hat F$ (that is, for each $n$-tuple $(t_1, \, ... \, , t_n) \in \{0,1\}^n$, $g (t_1, \, ... \, , t_n) = \hat F (t_1, \, ... \, , t_n)$).
	\end{theorem}

	\begin{proof}
		Let $g: \{0,1\}^n \rightarrow \{0,1\}$, let us build $F \in \fln$ such that $g = \hat F$. We can build this truth table:
		$$
		\begin{array}{c c c c | c}
			t_1 & t_2 & \cdots & t_n & g (t_1, \, ... \, , t_n) \\
			\hline
			0 & 0 & \cdots & 0 & b_0 \\
			0 & 0 & \cdots & 1 & b_1 \\
			&& \vdots && \vdots \\
			a_{i1} & a_{i2} & \cdots & a_{in} & b_i \\
			&& \vdots && \vdots \\
			1 & 1 & \cdots & 1 & b_{2^{n}-1} \\
		\end{array}
		$$

		\begin{osservation*}
		If $g(t_1, \, ... \, , t_n) = 0$ for all $(t_1, \, ... \, , t_n) \in \{0,1\}^n$ ($g$ always $0$), we can take $\bot = F$
		$$ \hat F = \hat \bot = 0 = g $$
		\end{osservation*}

		\noindent
		Otherwise, at least a bit is 1, let $I \subseteq \{0,1, \, ... \, , 2^{n}-1\}$ be defined by $i \in I$ \textbf{iff} $b_i = 1$, by our earlier observation $I \neq \emptyset$.

		\noindent
		For all $i \in I$ let us consider the $i$-th row of the table:
		$$ a_{i1} \ \ a_{i2} \ \ \cdots \ \ a_{in} \ \ | \ \ b_i = 1$$ \marginnote{$\leftarrow$ ($b_i =1$ since $i \in I$.)}[-0.9cm]
		\vspace{-0.75cm}

		\noindent
		For each $j \in \{1, \, ... \, , n\}$ we set
		$$
		\begin{cases}
			A_{ij}:= p_j & \text{ if } a_{ij} = 1 \\
			A_{ij}:= \neg p_j & \text{ if } a_{ij} = 0 \\
		\end{cases}
		$$

		\noindent
		We set \marginnote{$\leftarrow$ (note $F \in \fln$.)}[0.9cm]
		$$ F := \bigvee_{i \in I} \bigwedge_{j=1}^n A_{ij} $$

		\noindent
		We now have to prove that $\hat F = g$.\\

		\noindent
		It's enough to prove that for each $i \in I$:
		\begin{equation*}
		  \begin{rcases}
			\widehat{\left(\bigwedge_{j=1}^n A_{ij}\right)} (t_1, \, ... \, , t_n) = 1 \\
			\quad\quad\quad\quad\quad\quad \textbf{iff} \\
			(t_1, \, ... \, , t_n) =  (a_{i1}, a_{i2}, \, ... \, , a_{in})
		  \end{rcases}
		  (\dag)
		\end{equation*}

		\noindent
		Proof of $(\dag)$: Clearly
		$$ \widehat{\left(A_{i1} \wedge A_{i2} \wedge \, ... \, \wedge A_{in}\right)} (a_{i1}, a_{i2}, \, ... \, , a_{in}) = 1 $$
		For convenience:
		$$ \left(A_{i1} \wedge A_{i2} \wedge \, ... \, \wedge A_{in}\right) = \bigwedge_{j = 1}^n A_{ij} $$

		\noindent
		Since, by definition of $A_{ij}$
		$$ \hat A_{ij} (a_{i1}, a_{i2}, \, ... \, , a_{in})  = \begin{cases}
			a_{ij} & \text{ if } a_{ij} = 1\\
			1 - a_{ij} & \text{ if } a_{ij} = 0
		\end{cases}$$
		we have that $\hat A_{ij} (a_{i1}, \, ... \, , a_{in}) = 1$, So.
		$$ \widehat {\left(\bigwedge_{j = 1}^n A_{ij} \right)} (a_{i1}, a_{i2}, \, ... \, , a_{in}) = 1$$
		By definition of $I_\wedge$.\\

		\noindent
		Let now us take $(t_1, \, ... \, , t_n) \in \{0,1\}^n$, with $(t_1, \, ... \, , t_n) \neq (a_{i1}, a_{i2}, \, ... \, , a_{in})$ (so $(t_1, \, ... \, , t_n)$ is not the $i$th row). Let $j \in \{1, \, ... \, , n\}$ be such that $t_j \neq a_{ij}$ (there must be at least one position in which they are different, otherwise they would be the same row)
		$$
		\begin{cases}
			t_j = 1 & \text{ iff } a_{ij} = 0 \\
			t_j = 0 & \text{ iff } a_{ij} = 1 \\
		\end{cases}
		$$
		then
		$$ \hat A_{ij} (t_1, \, ... \, , t_n) =
		\begin{cases}
			t_j = 1 - a_{ij} = 0 & \text{ iff } t_j = 0 \text{ iff }  a_{ij} = 1 \\
			1 - t_j = a_{ij} = 0 & \text{ iff } t_j = 1 \text{ iff }  a_{ij} = 0
		\end{cases}
		$$
		So
		$$ \widehat{\left( \bigwedge_{j=1}^n A_{ij} \right)}(t_1, \, ... \, , t_n) = 0$$
		By definition of $I_\wedge$ . So we proved $(\dag)$.\\
		\noindent
		So, by $(\dag)$
		$$ \widehat{\left( \bigwedge_{j=1}^n A_{ij} \right)} (t_1, \, ... \, , t_n) = 1 \; \text{ iff } \; t_1, \, ... \, , t_n \; \text{ is the } i \text{th row} $$
		So \marginnote{$\leftarrow$ (by definition of $I$)}[1.4cm]
		$$\widehat{\left( \bigvee_{i \in I} \bigwedge_{j=1}^n A_{ij}\right)} (t_1, \, ... \, , t_n) = 1 \; \text{\textbf{iff}\; $(t_1, \, ... \, , t_n)$ is the $i$th row for some $i \in I$ \textbf{iff} $g(t_1, \, ... \, , t_n) = 1$.\\}$$

		\noindent
		So we have concluded that $F$, defined as:
		$$ F := \bigvee_{i \in I} \bigwedge_{j=1}^n A_{ij}$$
		is such that
		$$ \hat F (t_1, \, ... \, , t_n) = g(t_1, \, ... \, t_n) $$
		for all $(t_1, \, ... \, t_n) \in \{0,1\}^n$, so
		$$ \hat F = g $$
	\end{proof}

	\noindent
	Exercise: prove that for any $f: \{\}^n \rightarrow \{0,1\}$ there is a formula $G \in \fln$ such that $\hat G = f$ and
	$$G = \bigwedge_{i \in J} \bigvee_{j=1}^n B_{ij} $$
	where $J \subseteq \{0, \, ... \, , 2^n - 1\}$ such that $i \in J$ iff $b_j = 0$ and
	$$
	\begin{cases}
		B_{ij} = \neg p_j & \text{ if } a_{ij} = 1 \\
		B_{ij} = p_j & \text{ if } a_{ij} = 0
	\end{cases}
	$$
	Dual of before, just switch everything.\\

	\begin{figure}[H]
	\centering
	\includegraphics[width=1\textwidth]{img2.png}
	\end{figure}

	\noindent
	\textbf{Observation}: The functional completeness theorem \textbf{does not hold} for the \textbf{set of all formulas} $\FL$. It \textbf{works only when we restrict} to formulas and functions on \textbf{the first} $n$ \textbf{propositional letters}, for any $n \in \mathbb{N}$
	$$ Func^{(n)} = \hfln \cong \flne $$ \marginnote{$\leftarrow$ (Functional completeness theorem proves the first $=$.)}[-1cm]
	\vspace{-0.5cm}

	\noindent
	If we consider all propositional letters (remember $|\mathcal{L}| = \aleph_0$)
	$$ Func^{(\omega)} = \{f: \{0,1\}^\omega \rightarrow \{0,1\}\}$$
	But $\{0,1\}^\omega$ is the set of all infinitely long sequences of bits.\\
	$$ |Func^{(\omega)}| = |[0,1] \cap \mathbb{R}| = |\mathbb{R}| = 2^{\aleph_0} = \aleph_1 > \aleph_0 = |\fle| $$

	\noindent
	The normal forms:
	\begin{itemize}
		\item \textbf{Conjunctive}
		$$ \bigwedge_{i \in J} \bigvee_{j=1}^n B_{ij} $$
		\item \textbf{Disjunctive}
		$$ \bigvee_{i \in I} \bigwedge_{j=1}^n A_{ij} $$
	\end{itemize}

	Are said \textbf{canonical} or \textbf{complete}, conjunctive/disjunctive normal forms, that is every one of the disjuncts/conjuncts \textbf{mentions all the letters} $p_1, p_2, \, ... \, , p_n$.

	\begin{figure}[H]
	\centering
	\includegraphics[width=1\textwidth]{img3.png}
	\end{figure}

	\section{Functionally complete set of connectives}

	A \textbf{set} $C \subseteq \{\wedge, \vee, \rightarrow, \neg, \leftrightarrow, \top, \bot, \, \dots \, \}$ (subset of connectives) is \textbf{functionally complete} iff for all $F \in \fln$ (for all $n \in \mathbb{N}$) there is some $G \in \fln$ written using only the connectives in $C$ such that $\hat F = \hat G$ (or $F \equiv G$, or $[F]_\equiv = [G]_\equiv$, or $\models F \leftrightarrow G$, they're all equivalent). \\

	\noindent
	Some examples of functional complete sets are:
	\vspace{-0.2cm}
	\begin{itemize}
		\item $\{\wedge, \vee, \neg \}$: is \textbf{functionally complete}, by the \textbf{functional completeness theorem}.
		\item $\{\neg, \wedge\}$: same as above, use De Morgan laws to obtain this set
		\item $\{\neg, \vee\}$: also De Morgan
		\item $\{\neg, \rightarrow\}$: can be proven remembering that $A \wedge B \equiv (\neg A) \rightarrow B$
		\item $\{\rightarrow, \bot\}$: since $\neg A \equiv A \rightarrow \bot$
	\end{itemize}

	\paragraph{NAND:} Operator defined as negation of AND, it's usually defined as:
	$$ A \nand B := \neg (A \wedge B) $$
	$$ I_{\nand} = \{0,1\}^2 \rightarrow \{0,1\} $$
	$$ I_{\nand} (x,y) = 1 - \min\{x,y\}$$
	$$
	\begin{array}{c c | c c | c}
		A & B & \neg & (A \wedge B) & \nand \\
		\hline
		0 & 0 & 1 & 0 & 1 \\
		0 & 1 & 1 & 0 & 1 \\
		1 & 0 & 1 & 0 & 1 \\
		1 & 1 & 0 & 1 & 0 \\
	\end{array}
	$$


	\noindent
	Recalling that $\{\neg, \wedge\}$ is functionally complete, the set only composed by \textbf{NAND} is \textbf{functionally complete}:
	$$
	\begin{cases}
		\neg A \equiv A \nand A \\
		A \wedge B \equiv \neg (A \nand B) \equiv (A \nand B) \nand (A \nand B)
	\end{cases}
	$$

	\marginnote{EXERCISES}
	\vspace{-0.5cm}

	\section{Normal Forms}

	We have seen canonical/complete normal forms (CNF, DNF) and they are \textbf{generally too long}.\\
	We want to produce CNFs and DNFs without their being canonical, in order to obtain shorter formulas.

	\subsection{Implication Free Normal Form IFNF:}

	We're getting rid of implication.

	\begin{definition} A Formula $F \in \FL$ is in IFNF ($F \in IFNF$) \textbf{iff} $F$ \textbf{does not contain} any occurrence \textbf{of} $\rightarrow$.
	\end{definition}

	\subsection{Negation Normal Form NNF}

	\begin{definition} A formula $F \in IFNF$ is in NNF ($F \in NNF$) iff the \textbf{negation connective} is applied \textbf{only to propositional letters} (if $\neg B$ is a sub-formula of $F$, then $B = p_i$ for some $p_i \in \mathcal{L}$).
	 \end{definition}

	\begin{example} Some Examples:
	  \begin{itemize}
	  \item Is $p_1 \wedge \neg (p_2 \vee p_3) \in $ NNF?
		Obviously no, but its equivalent $p_1 \wedge (\neg p_2 \wedge \neg p_3)$
	  \item Is $p_1 \rightarrow \neg p_2 \in$ NNF? No, since it's not $\in IFNF$ (the negation is only on letters though)
	  \end{itemize}
	\end{example}

	\noindent
	Some interesting relations are:
	$$
	\begin{array}{c c}
		CNF & \subseteq \\
		DNF & \subseteq
	\end{array}
	NNF \subseteq IFNF
	$$

	%Check below

	\paragraph{Notation:} $G \preceq F$, $G$ is a \textbf{Subformula} of $F$ (for $G,F \in \FL$) if $G$ occurs in every $L$-construction of $F$ (you need $G$ to build $F$, get the minimal $L$-construction and check if $G$ is in there).\\

	\paragraph{Transformations:} We shall use the following equivalences to \textbf{replace subformulas} with \textbf{equivalent ones}:
	\begin{enumerate}
		\item $C \rightarrow D \equiv (\neg C) \vee D$
		\item $\neg \neg C \equiv C$
		\item $\neg (C \vee D) \equiv \neg C \wedge \neg D$ (De Morgan)
		\item $\neg (C \wedge D) \equiv \neg C \vee \neg D$ (De Morgan)
	\end{enumerate}

	We shall use these logical equivalences to transform subformulas of a given formula $F \in \FL$ using these from left to right, we're \textbf{using them as transformation rules} (we could do the opposite, but goes against our normal forms).

	\begin{lemma}
	For any $F \in \FL$ we can build another formula $G \in \FL$ (in finite number of steps) such that $F \equiv G$ and $G \in IFNF$. Any formula can be "transformed" into an equivalent one in IFNF.
	\end{lemma}
	\begin{proof}
		For any $E \in \FL$ let $i(E)$ be the number of occurrences of $\rightarrow$ in $E$. We can observe that $E \in IFNF$ iff $i(E) = 0$. If $E \notin IFNF$ then $i(E) > 0$.

	\begin{enumerate}
		\item Take $C \rightarrow D \preceq E$ (which exists for any $E$ with $i(E) > 0$).

		\item Replace $C \rightarrow D$ with $(\neg C) \vee D$ and call $E'$ the formula obtained.
		$$ E' = E\left[\sfrac{(\neg C) \vee D}{C \rightarrow D}\right]$$
		(with a small abuse of notation, I think we defined the substitution only for letters but I think it's admissible).
	\end{enumerate}

		\noindent
		Now $E \equiv E'$ for the transformation at step $(2)$ thanks to the substitution theorem and now we have
		$$ i(E) > i(E')$$
		since we removed at least one implication (possibly more than one).

		So, given in input $F \in \FL$ with $F \notin IFNF$, we build a sequence
		$$ F = F_0 \equiv F_1\equiv \, \dots \, \equiv F_u \in IFNF \;\;\;\;\;\;\;\; (\dag) $$
		where
		$$ i(F_0) > i(F_1) > \, \dots \, > i(F_n) = 0 $$
		Since we're strictly descending natural numbers since each $F_{i+1}$ is obtained from $F_i$ by substituting some $C \rightarrow D \preceq F$ with $(\neg C) \vee D$. Trivially, the sequence $(\dag)$ ends in a finite number of steps $u$ with $u \leq i(F_0) = i(F)$, producing a formula $F_u$ such that $F_u \equiv F$, $F_u \in IFNF$.
	\end{proof}

	\begin{lemma} for any $F \in IFNF$ we can build, in a finite number of steps, a \textbf{sequence}
	$$ F = F_0 \equiv F_1 \equiv \, \dots \, \equiv F_v = G $$
	such that $G \equiv F$ and $G \in NNF$.
	\end{lemma}
	\begin{proof}

	For any $A \in IFNF$ let $\ell(A) :=$ the \textbf{number of} occurrences of \textbf{connectives} $\{\wedge, \vee, \neg\}$ in $A$.

	\newpage

	\noindent
	For any $E \in IFNF$ let
	$$m (E):= \sum_{\neg A \preceq E} \ell(A)$$

	\noindent
	Lets prove that $m(E) = 0$ \textbf{iff} $E \in NNF$.

	  \begin{itemize}
		\item ($\Rightarrow$) If $m(E) = 0$ then by definition we have $\sum_{\neg A \preceq E} 0$ then $\ell(A) = 0$ for all $\neg A \preceq E$, then $A = p$ for some $p \in \mathcal{L}$, then $E \in NNF$ (by definition of NNF).

		\item ($\Leftarrow$) If $E \in NNF$ then $\neg A \preceq E$ iff $A = p$ for some $p \in \mathcal{L}$, so if $A$ is a propositional letter, it contains no connectives $\ell (A) = 0$, so $m(E) = \sum_{\neg A \preceq E} 0 = 0$.
	  \end{itemize}

	We now have to show that each application of our transformations (we list here) is such that one step from $F_i$ to $F_{i+1}$ has $m(F_i) > m(F_{i+1})$.\\

	\noindent
	Cases:
	\begin{itemize}
		\item \textbf{Transformation}: $\neg \neg C \leadsto C$ \\
		Using $\neg \neg C \preceq F_i$ we need to prove
		$$ m(\neg \neg C) > m(C) $$

		By definition:
		$$m(C) = \sum_{\neg A \preceq C} \ell(A)$$

		So:
		\begin{align*}
		  m(\neg \neg C) = \sum_{\neg A \preceq \neg \neg C} \ell(A) &= \ell (\neg C) + \sum_{\neg A \preceq \neg C} \ell (A) \\
		  &= \ell (C) + 1 + \ell (C) + \sum_{\neg A \preceq C} \ell (A)
		\end{align*}

		So $m(\neg \neg C) > m(C)$ by at least $1$.\\

	  \item \textbf{Transformation}: $\neg (C \wedge D) \leadsto (\neg C \vee \neg D)$ \\
		Using $\neg (C \wedge D) \preceq F_i$ we need to prove:
		$$ m(\neg (C \wedge D)) > m(\neg C \vee \neg D) $$

		By definition:
		$$m(\neg C \vee \neg D) = \sum_{\neg A \preceq \neg C \vee \neg D} \ell (A) = \ell (C) + \ell(D) + \sum_{\neg A \preceq C} \ell(A) + \sum_{\neg A \preceq D} \ell(A)$$

		So:
		\begin{align*}
		  m(\neg (C \wedge D)) = \sum_{\neg A \preceq \neg (C \wedge D)} \ell(A) &= \ell (C \wedge D) &+ \sum_{\neg A \preceq C} \ell(A) + \sum_{\neg A \preceq D} \ell(A)  \\
		  &= \ell (C) + 1 + \ell (D) &+ \sum_{\neg A \preceq C} \ell(A) + \sum_{\neg A \preceq D} \ell(A)
		\end{align*}

		then $m(\neg (C \wedge D)) > m(\neg C \vee \neg D)$.\\

		\item \textbf{Transformation}: $\neg (C \vee D) \leadsto (\neg C \wedge \neg D)$ \\
		  Using $\neg (C \vee D) \preceq F_i$ we need to prove:
		  $$ m(\neg (C \vee D)) > m(\neg C \wedge \neg D) $$
		  Same as before, just switch symbols.

	\end{itemize}
	\end{proof}

	\noindent
	\textbf{Observe}: if $G \preceq F$, going from $G$ to $H$ and subsequently from $F$ to $F'$, $F' = F[\sfrac{H}{G}]$
	$$ \text{if } \; m (G) > m(H) \; \text{ then } \; m(F) > m(F')$$

	\noindent
	As a matter of fact:
	$$m(F) = \sum_{\neg A \preceq F} \ell(A) = \sum_{\neg A \preceq G} \ell(A) + \sum_{\neg A \preceq F, \neg A \not \preceq G} \ell (A)  = m(G) + \sum_{\neg A \preceq F, \neg A \not \preceq G} \ell (A) $$
	$$ m(F') = \sum_{\neg A \preceq F'} \ell(A) = \sum_{\neg A \preceq H} \ell(A) + \sum_{\neg A \preceq F, \neg A \not \preceq G} \ell (A)  = m(H) + \sum_{\neg A \preceq F, \neg A \not \preceq G} \ell (A)$$
	$$ \implies m(F) > m(F') $$

	\begin{figure}[H]
	\centering
	\includegraphics[width=1\textwidth]{img4.png}
	\end{figure}

	\vspace{-1.5cm}

	\begin{figure}[H]
	\centering
	\includegraphics[width=1\textwidth]{img5.png}
	\end{figure}

	\subsection{Conjunctive Normal Form CNF}

	\paragraph{Terminology: }
	\begin{itemize}
		\item \textbf{Definition: Literal}, by literal we mean a formula in the form $p$ or $\neg p$ for some $p \in \mathbb{L}$. A letter or a negation of one.\\

		\item \textbf{Definition: Opposite Literal}, if $A$ is a literal, then we denote by $\overline{A}$, called opposite of $A$, the literal so defined:
		$$
		\begin{cases}
			\text{if } \; A = p & \text{ then } \overline{A} = \neg p \\
			\text{if } \; A = \neg p & \text{ then } \overline{A} = p
		\end{cases}
		$$
		Equivalently
		$$ \overline{A} := (\neg A)^N$$

		\item \textbf{Definition: Clause}, a clause is a disjunction of a finite number of literals:
		$$ l_1 \vee l_2 \vee \, \dots \, \vee l_k $$
		where each $l_i$ is a literal.\\

		\item \textbf{Definition: CNF}, a CNF is a conjunction of a finite number of clauses:
		$$ C_1 \wedge C_2 \wedge \, \dots \, \wedge C_n$$
		where each $C_j$ is a clause.\\

		\item \textbf{Definition: Empty Clause}, an empty clause ($\square$) is, by definition, a disjunction of zero literals.\\

		\item \textbf{Definition: Empty CNF}, the empty CNF ($\emptyset$) is the conjunction of zero clauses.
	\end{itemize}

	\noindent
	Thinking about $l_1$ as a literal, and the disjunctions

	\begin{align*}
	  	&0 \ \ \square \\
		&1 \ \ l_1 \\
		&2 \ \ l_1 \vee l_2 \\
		&3 \ \ l_1 \vee l_2 \vee l_3 \\
	  	&\quad \quad \quad \vdots \\
	  	&k \ \ l_1 \vee l_2 \vee l_3 \vee \, \dots \, l_k \\
	  	&\quad \quad \quad \vdots \\
	\end{align*}

	If any of these is satisfiable, then all of the ones below are satisfiable, so it makes sense to put $\square$ on top of the first literal $l_1$, so, by definition, the empty literal $\square$ \textbf{is unsatisfiable and a contradiction}.
	Same game with clauses and empty CNF $\emptyset$:
	\begin{align*}
	  	&0 \ \ \emptyset \\
		&1 \ \ C_1 \\
		&2 \ \ C_1 \wedge C_2 \\
	  	&\quad \quad \quad \vdots \\
	\end{align*}

	But here the satisfiability flows upwards, if one is satisfied then all of the above are satisfied. The empty CNF $\emptyset$ is still above all, and as such is satisfied by all assignments and is a tautology

	\subsection{Disjunctive Normal Form DNF}

	A DNF is a disjunction of a finite number of formulas:
	$$ D_1 \vee D_2 \vee \, \dots \, \vee D_h $$
	Where each $D_i$ (dual of a clause, hasn't got a proper name, maybe co-clause) is a conjunction of a finite number of literals:
	$$ l_1 \wedge l_2 \wedge \, \dots \, \wedge l_m $$

	\noindent
	Every definition stated before has a DNF dual but they don't have a proper name since \textit{nobody cares}.

	\begin{lemma} Every formula $F \in \FL$ is logically equivalent to a formula $F^C \in CNF$ and to a formula $F^D \in DNF$
	$$ (CNF \ni F^C \equiv F \equiv F^D \in DNF )$$
	\end{lemma}

	\begin{proof}
		By structural induction of $F$:
		\begin{itemize}
			\item \textbf{Base}: $F = p$, $p \in \mathcal{L}$ then $p^C \equiv p \equiv p^D$. A single literal is already in CNF and DNF, $p^C = p = p_D$.\\

			\item \textbf{Step}:
			\begin{itemize}
			\item Case $F = \neg G$, by I.H., there are $G^c \in CNF$, $G^D \in DNF$ with $G^C \equiv G \equiv G^D$
			$$ F = \neg G \equiv \neg(G^D) \in CNF $$
			$$ F = \neg G $$
			$$ \neg (G^D) = \neg (D_1 \vee D_2 \vee \, \dots \, \vee D_v) $$

			with each $D_i$
			$$ D_i = l_{i1} \wedge l_{i2} \wedge \, \dots \, \wedge l_{iq}$$

			but by De Morgan
			$$ \neg (G^D) \equiv \neg D_1 \wedge \neg D_2 \wedge \, \dots \, \wedge \neg D_v $$

			And, still by De Morgan, each $\neg D_i$
			$$ D_i \equiv (\overline{l}_{i1} \vee \overline{l}_{i2} \vee \, \dots \, \vee \overline{l}_{iu_i}) $$

		So $\neg (G^D)$ is the conjunction $\neg D_1 \neg D_2 \wedge \, \dots \, \neg D_v$ of a finite number of formulas $(\neg D_i)$, each one equivalent to a disjunction of a finite number of literals, that is $\neg (G^D) \in CNF$, and $F \equiv \neg G \equiv \neg (G^D)$.\\
			\end{itemize}

		\begin{itemize}
		\item Case: $F = (G \wedge H)$, by I.H.
		$$  G^C, H^C \in CNF, \;\;\;\; G^D, H^D \in DNF $$
		$$ \text{with } \; G^C \equiv G \equiv G^D, \;\; H^C \equiv H \equiv H^D $$

		We want to build \textbf{2 new formulas}:
			\item $F \equiv F^D$, defining $F^D$ as  $:= G^D \vee H^D \in DNF$, then by definition is $ \in DNF$, and by construction $G^D \vee H^D \equiv F$. Disjunction of disjunctions of finitely many co-clauses, so still $\in DNF$.\\

			\item $F \equiv F^C$, then
			$$ G^C = G_1 \wedge G_2 \wedge \, \dots \, \wedge G_u $$
			$$ H^C = H_1 \wedge H_2 \wedge \, \dots \, \wedge H_v $$
			So
			$$ G \vee H \equiv G^C \vee H^C \equiv (G_1 \wedge G_2 \wedge \, \dots \, \wedge G_u) \vee (H_1 \wedge H_2 \wedge \, \dots \, \wedge H_v) \equiv $$
			Using generalized distributivity:
			$$ \equiv \bigwedge_{i=1}^u \bigwedge_{j=1}^v (G_i \wedge H_j) $$

		Each $G_i$ and each $H_j$ are clauses, i.e., disjunctions of finitely many literals, so each term in the form $G_i \vee H_j$ is again a clause, and then the whole formula is a finite conjunction of clauses and, by definition, a CNF, also $\equiv G \vee H = F$. \\
			\item 	Then we have the case $(G \wedge H)$, which is left as an exercise to the reader. \\
			\item  Case (to complete our set of connectives): $(G \rightarrow H)$, which is also left as an exercise, because implication can be reduced to $G \rightarrow H \equiv (\neg G) \vee H$.\\
		\end{itemize}
	\end{itemize}
	\end{proof}

	\begin{figure}[H]
	\centering
	\includegraphics[width=1\textwidth]{img6.png}
	\end{figure}

	\newpage

%%% Local Variables:
%%% TeX-master: "../MathLogic.tex"
%%% End:

